The evaluation protocol changes which aggregation claim the evidence supports.
Independent paired simulations establish a negative/positive state-ranking
contrast at two selected duration configurations, with switch coverage changing
too. The symmetric-generator study reproduces that contrast at the original
window setting; point-estimate exceptions elsewhere have intervals crossing
zero and do not establish population sign reversals. These are configuration
comparisons, not an isolated causal duration law.
Non-clustering shrinkage can improve pooled ranking with temporal order intact,
so structural gain alone does not identify a state-recovery mechanism.
The residual decomposition locates a substantial gain in restoring comparisons
that hard emission ties, rather than in better state assignments.
Raw fallback isolates loss from constant rejection output while leaving gate
selection as a prototype-specific issue. The PELT plus VAR comparator shows
that pooled improvement without an established state-ranking gain also occurs
in a different segmentation/refit structure.

Boundary judgments require explicit scale and zero-score conventions. The
common raw-amplitude F1 decreases do not by themselves establish poorer
boundary localisation: method-relative thresholds and standard AP change those
assessments. Both conventions, fixed threshold curves and stable-edge rates
therefore belong alongside temporal state AUROC. The ranking principles are
established, and the finite identities are diagnostic; their contribution here
is an executable control protocol and evidence about which interpretations
survive its distinctions in the tested linear systems.
